% Chapter 8 supplemental resources. Included by ch08_resources.tex.
\section{Supplemental Resources}
\label{sec:comp-ch08-supplemental}

These supplemental resources translate the Chapter~8 concepts into editable
institutional records. They should be adapted to the intended use and consequence,
completed by a multidisciplinary team, and maintained under version control. They
are not validated clinical, legal, regulatory, procurement, accreditation, or health
technology assessment instruments.

\begin{longtable}{@{}p{0.28\textwidth}p{0.38\textwidth}p{0.26\textwidth}@{}}
\caption{Chapter 8 supplemental-resource set}
\label{tab:comp-ch08-supplemental-set}\\
\toprule
\textbf{Resource} & \textbf{Purpose and minimum content} & \textbf{Expected output} \\
\midrule
\endfirsthead
\toprule
\textbf{Resource} & \textbf{Purpose and minimum content} & \textbf{Expected output} \\
\midrule
\endhead
\bottomrule
\endlastfoot
Model Specification & Define scope, owner, evidence source, assumptions, decision threshold, review date, and required action. & Blank record, synthetic example, and version history. \\
Verification, Validation, and Reproducibility & Define scope, owner, evidence source, assumptions, decision threshold, review date, and required action. & Blank record, synthetic example, and version history. \\
Uncertainty and Value-of-Information Record & Define scope, owner, evidence source, assumptions, decision threshold, review date, and required action. & Blank record, synthetic example, and version history. \\
\end{longtable}

\subsection{Model Specification}
\label{sec:comp-ch08-supp-1}
Use this record to establish the decision before analysis begins. Complete every
field or explain why it is not applicable. Do not replace missing evidence with an
implicit assumption.

\begin{longtable}{@{}p{0.27\textwidth}p{0.65\textwidth}@{}}
\caption{Minimum specification fields}
\label{tab:comp-ch08-spec-fields}\\
\toprule
\textbf{Field} & \textbf{Required entry} \\
\midrule
\endfirsthead
\toprule
\textbf{Field} & \textbf{Required entry} \\
\midrule
\endhead
\bottomrule
\endlastfoot
Decision and authority & Decision to be made; accountable sponsor; decision owner; approval and escalation authority. \
Population and setting & Included and excluded populations, sites, services, users, and time period. \
Problem and baseline & Current pathway, burden, variation, unmet need, and baseline measures. \
Comparator set & Current practice and other feasible clinical, staffing, process, digital, shared-service, and no-change alternatives. \
Intervention boundary & Data, model or method, interface, people, workflow, capacity, implementation, monitoring, vendor support, and exit. \
Evidence and uncertainty & Sources, dates, versions, applicability, confidence, gaps, and assumptions. \
Costs and consequences & Perspective, budget holder, resource use, outcomes, workload, distribution, and opportunity cost. \
Safeguards & Safety, lawful use, privacy, cybersecurity, equity, meaningful oversight, continuity, and contestability. \
Decision conditions & Proceed, redesign, restrict, pilot, defer, replace, or retire criteria; review date; trigger events. \
\end{longtable}

\subsection{Verification, Validation, and Reproducibility}
\label{sec:comp-ch08-supp-2}
The following sequence converts findings into accountable action. The local record
should name the responsible person for each activity, the completion date, the
evidence of completion, and any residual uncertainty.

\begin{enumerate}
 \item Specify the comparator, human reading workflow, priority threshold, queue discipline, and downstream pathway..
 \item Use a classification decision tree to estimate true-positive, false-negative, false-positive, and true-negative consequences..
 \item Use discrete-event simulation or a queueing model to estimate time-to-report distributions under constrained reader capacity..
 \item Link material diagnostic delay to downstream cost and QALY consequences..
 \item Represent subgroup prevalence, performance, access, and delay distributions explicitly..
 \item Conduct deterministic sensitivity analysis, probabilistic sensitivity analysis, and value-of-information analysis..
\end{enumerate}

\subsection{Uncertainty and Value-of-Information Record}
\label{sec:comp-ch08-supp-3}
A final record should distinguish observations, modelled estimates, value
judgements, and contractual claims. It should preserve the evidence cut-off, exact
system or instrument version, dissent, unresolved risks, monitoring plan, and the
reason the selected action is preferred to the alternatives.

\begin{longtable}{@{}p{0.24\textwidth}p{0.31\textwidth}p{0.29\textwidth}@{}}
\caption{Decision, monitoring, and review record}
\label{tab:comp-ch08-decision-record}\\
\toprule
\textbf{Record element} & \textbf{Required content} & \textbf{Review evidence} \\
\midrule
\endfirsthead
\toprule
\textbf{Record element} & \textbf{Required content} & \textbf{Review evidence} \\
\midrule
\endhead
\bottomrule
\endlastfoot
Decision & Selected action and rejected alternatives, with reasons. & Signed and dated decision; declared conflicts and dissent. \
Conditions & Preconditions, restrictions, milestones, and contractual duties. & Completion evidence and accountable owner. \
Indicators & Leading, process, outcome, balancing, subgroup, cost, and governance measures. & Data source, definition, threshold, frequency, and owner. \
Triggers & Events requiring investigation, revalidation, restriction, or exit. & Incident, material-change, or threshold-breach record. \
Review & Scheduled date and event-triggered review route. & Current evidence, realized value, residual risk, and alternatives. \
Closure & Transition, communication, data closure, recovery, and lessons. & Verified closure record and organizational-learning entry. \
\end{longtable}
